\documentclass{article}

\usepackage{neurips_2026}

\usepackage[utf8]{inputenc}
\usepackage[T1]{fontenc}
\usepackage{hyperref}
\usepackage{url}
\usepackage{booktabs}
\usepackage{amsfonts}
\usepackage{amsmath}
\usepackage{nicefrac}
\usepackage{microtype}
\usepackage{xcolor}
\usepackage{graphicx}

% Anything still open. Delete the macro before submission: if a \TBD survives to
% the deadline, that claim comes out of the paper.
\newcommand{\TBD}[1]{\textcolor{red}{\textbf{[TBD: #1]}}}
% Space held for a figure that is not drawn yet, with what it has to show.
\newcommand{\figstub}[2]{%
  \fbox{\begin{minipage}[c][#1][c]{0.97\linewidth}%
    \centering\small\color{red}\textbf{[FIGURE]}\\[2pt]\itshape #2%
  \end{minipage}}}

\title{\TBD{title. Working: Early Warning Is Not a Property of the Probe}}

\author{%
  \TBD{authors, anonymised for submission}
}

\begin{document}

\maketitle

\begin{abstract}
\TBD{written last, once the storyline is fixed. It has to state the population
and its natural degeneration rate, the separation of detection from pre-onset
warning, the size of the pre-onset signal against a causal baseline, the
decomposition into position, length, prompt and domain, and the held-out
confirmation of the domain result.}
\end{abstract}

\section{Introduction}
\label{sec:intro}

\TBD{written last. The order of the argument is settled even if the framing is
not, and it is: (i) a loop that has started is easy to see and the metric
everyone reports is saturated, so it cannot rank methods or even choose a layer;
(ii) the question that matters is whether anything is readable \emph{before} the
loop, which needs tokens scored by position relative to an onset; (iii) an
apparent early warning has cheaper explanations than an impending loop, and this
corpus is built so that each of them is measurable; (iv) what survives is small,
concentrated in one kind of text, and absent late in an answer.}

\paragraph{Contributions.}
\begin{enumerate}
  \item A measurement that separates recognising a loop in progress from seeing
  it coming, by scoring every token according to its position relative to a
  semantically located onset, at operating points fixed by a per-answer
  false-alarm budget (Section~\ref{sec:protocol}).
  \item A decomposition of the apparent pre-onset signal into absolute position,
  answer length, prompt-level predisposition and domain, each against a matched
  null, with the residual reported rather than assumed (Section~\ref{sec:controls}).
  \item The observation that the pre-onset signal is concentrated in one kind of
  text rather than general, confirmed on a domain held out of training entirely
  (Section~\ref{sec:exp-domain}).
  \item A bridge between this protocol and the reporting conventions of the
  nearest prior work, holding the scorer fixed, which quantifies how much of the
  disagreement in published early-warning numbers is a matter of convention
  (Section~\ref{sec:exp-bridge}).
  \item The finding that answer-level detection is saturated across almost the
  whole depth of the network while pre-onset coverage varies by more than an
  order of magnitude over the same depths, so the metric this literature reports
  cannot see the choice that determines early warning (Section~\ref{sec:exp-depth}).
\end{enumerate}

\section{Related work}
\label{sec:related}

\paragraph{Repetitive degeneration.}
Repetition is a long-standing failure of neural text generation. It has been
attributed to a positive feedback loop in which the probability of a phrase
rises with each repetition \citep{holtzman2020curious}, to high-inflow states in
the transition dynamics \citep{fu2021theoretical}, to a self-reinforcement effect
measured directly on sentence probabilities \citep{xu2022learning}, and to
repetition in the training data \citep{li2023repetition}. It persists in current
reasoning models, where it consumes the decoding budget
\citep{xie2025wordsalad,dong2025rethinking}. This work does not propose a
mechanism. It asks how much of the phenomenon is visible to a linear readout
before it appears in the text.

\paragraph{Measuring repetition.}
The standard measures are computed on a finished sequence: \texttt{distinct-n}
\citep{li2016diversity} and its complement \texttt{seq-rep-n}
\citep{welleck2020neural} count duplicate $n$-grams in a completed generation.
Sliding a fixed window along the sequence, as the moving-average type-token ratio
does \citep{covington2010cutting}, removes the confound between such a ratio and
the length of the sample. \citet{welleck2020neural} also define \texttt{rep/l},
the share of tokens already occurring in the previous $l$ tokens, which is the
literature's per-token repetition signal and is backward looking by
construction. Section~\ref{sec:baselines} uses both a windowed
\texttt{seq-rep-2} and \texttt{rep/l}.

\paragraph{Reading a loop from internal states.}
Loops are visible in hidden states once they have begun.
\citet{duan2026circular} report accuracy and area under the curve above $0.99$
separating loop from non-loop states with linear and non-linear classifiers on
averaged final-layer representations. \citet{xie2025wordsalad} train a linear
classifier on the hidden state of the token ending each reasoning chunk and
report roughly $90$ to $94\%$ accuracy at temperature zero. \citet{yu2025breaking}
detect adversarially induced loops from neuron activation patterns at $95.24\%$
accuracy. We take these results as settled and reproduce them. Our argument is
that the task they measure is close to saturated because of the shape of the
data, so a high answer-level number says little about whether anything is
readable before the loop begins. \TBD{verify directly that
\citet{duan2026circular} and \citet{xie2025wordsalad} both read the final layer.
This sentence is load-bearing for Section~\ref{sec:exp-depth} and currently rests
on a secondary reading.}

\paragraph{Predicting a loop.}
The closest work is \citet{duan2026circular}, who accumulate classifier evidence
with a cumulative-sum change-point rule \citep{page1954continuous} and report
early detection rates and lead times in tokens. Our results do not contradict
theirs; they are measured under different conditions, and those conditions are
what this paper is about. Their prediction experiments use greedy decoding,
balanced evaluation sets, sentence-level representations, and an operating point
far looser than one a serving system would accept. We sample, preserve the
natural base rate, score every token, and fix a budget on healthy answers.
Section~\ref{sec:exp-bridge} measures what each of those choices is worth.
\citet{duan2026circular} also report that semantic circularity precedes verbatim
repetition. We do not dispute that a precursor exists; we ask the narrower
question of whether it is linearly decodable per token at a false-alarm rate a
deployment would accept.

\paragraph{Detection strength grows inside a loop.}
\citet{xie2025wordsalad} frame their result as models being self-aware when
trapped in repetition. Their own appendix supports a more careful reading: for a
tracked example the classifier scores the first repetitive chunks at essentially
zero and reaches certainty only hundreds of chunks later. Being self-aware once
trapped and being self-aware on the approach are two claims, and the protocol of
Section~\ref{sec:protocol} is built to report them apart.

\paragraph{What a probe measures.}
A probe can reach high accuracy on labels that carry no information, so accuracy
alone does not establish that a representation contains the target
\citep{hewitt2019designing}. \citet{sahoo2026linear} give a recent example in
which probes at $100\%$ accuracy fall to chance once source identity and
response length are residualised out, which is directly relevant here because
every degenerate answer in our corpus reaches the token cap while healthy answers
are shorter. \citet{fomin2026internal} draw the distinction between an internal
readout that describes the current situation and one that predicts what happens
next, and report negative results for the second across three probe families.
That distinction is the one we draw between in-pattern and warning coverage; our
contribution is its instantiation in a setting where the onset can be located, so
the gap can be measured rather than argued. \citet{kramar2026building} report
that probes degrade under production distribution shift, particularly on long
contexts, for which Section~\ref{sec:controls} supplies a concrete mechanism.

\paragraph{Thresholded repetition rules in practice.}
Rules that threshold a repetition statistic already decide what a corpus keeps
\citep{rae2021scaling} and when a serving stack stops generating, and recent work
uses one to label loops and to trigger an intervention \citep{xu2026loopguard}.
None of them reports a detection rate or a false-alarm rate against an
independent reference. Appendix~\TBD{ref} measures three of these statistics
against a semantic label on the population where they are least trustworthy.

\section{Corpus and labels}
\label{sec:corpus}

\subsection{Generation}
The corpus is \TBD{36{,}300} answers generated from \texttt{Apertus-8B-Instruct}
\citep{hernandez2025apertus}, an openly released 8B instruction-tuned model, by sampling
$K=10$ continuations for each of \TBD{3{,}630} prompts at temperature $0.7$ and
top-$p$ $0.9$, with a hard cap of $4096$ tokens and a fixed seed. Prompts are
drawn from seven sources spanning mathematics, competitive programming, general
instruction following and medical reasoning. Five are used in domain and split by
prompt into training, validation and an in-domain test set; two are held out
entirely and read only for zero-shot evaluation.

Two properties of this design are used throughout. Sampling rather than greedy
decoding matches the generation policy a served model would use, and it produces
several trajectories per prompt, which is what makes prompt-level effects
measurable at all (Section~\ref{sec:controls}). Splits are assigned at the prompt
level, so all $K$ answers to a prompt share a split.

\subsection{What counts as degeneration}
An answer that ended at an end-of-sequence token is a \textbf{negative}. An
answer cut off by the token cap is sent to a language-model judge, which reads
the prompt and the answer and decides whether the text stopped making progress.
The prompt is essential: an enumeration, a code template, a brute-force search or
repetition the prompt explicitly requested are all legitimately repetitive, and a
purely structural rule cannot tell them from a loop.

Treating every naturally terminated answer as a negative is an assumption, and
Appendix~\TBD{ref} reports the evidence for it. \TBD{state the resulting scope
plainly: this corpus studies persistent degeneration that runs to the cap, not
every transient repetitive episode.}

\subsection{The degeneration frontier}
\label{sec:frontier}
When the judge finds degeneration it returns a short verbatim quote marking where
the pattern first begins. Asking for a quote rather than a position is what makes
the label checkable: a separate step locates the quote in the answer's own token
stream and records the index of its first token, giving the \textbf{frontier}
$f_r$ for answer $r$. A quote that cannot be located leaves its answer excluded
from every split, with the reason recorded, rather than given an invented onset.
\TBD{report the verbatim rate, the server-side re-verification, the refusal rate
and where it concentrates.}

Of the \TBD{36{,}300} answers, \TBD{890} reached the cap and \TBD{818} yield a
usable frontier, a natural degeneration rate of \TBD{2.3\%}. Table~\ref{tab:corpus}
gives the split populations.

The frontier is treated as a reference point, not as ground truth. It is
uncertain, and Section~\ref{sec:controls} reports how far the conclusions move
when it is displaced. \TBD{human validation of a stratified sample, with
inter-annotator agreement reported beside judge agreement so that judge error can
be read against task difficulty.}

\begin{table}[t]
  \caption{Corpus populations. A positive is an answer that reached the token cap
  and for which a frontier could be located. \TBD{fill}}
  \label{tab:corpus}
  \centering
  \begin{tabular}{lrrrr}
    \toprule
    Split & Prompts & Answers & Positives & Rate \\
    \midrule
    train                & \TBD{} & \TBD{} & \TBD{} & \TBD{} \\
    validation           & \TBD{} & \TBD{} & \TBD{} & \TBD{} \\
    test, in domain      & \TBD{} & \TBD{} & \TBD{} & \TBD{} \\
    test, held-out domains & \TBD{} & \TBD{} & \TBD{} & \TBD{} \\
    \bottomrule
  \end{tabular}
\end{table}

\subsection{Where loops start, and why it shapes everything}
\label{sec:onset-shape}
The single most consequential fact about this corpus is where loops begin. The
quartiles of the frontier position are \TBD{399}, \TBD{721} and \TBD{1{,}144}
tokens, and \TBD{about 3\%} of positives begin at the first token. Since every
positive runs to the cap, roughly \TBD{four fifths} of a positive answer is
already inside the loop. Any measure that lets those tokens dominate reads as
nearly perfect for almost any scorer, which is the reason for
Section~\ref{sec:protocol}.

\section{Probes and training}
\label{sec:probes}

\subsection{The probe}
A probe is a learned normalisation followed by one linear map, reading the
residual stream at a single layer at a single token and emitting a probability:
\begin{equation}
  p_t \;=\; \sigma\!\left(w^{\top}\mathrm{LN}\!\left(h^{(\ell)}_t\right) + b\right),
\end{equation}
about \TBD{12{,}289} parameters per depth. The score at a token is a function of
that token's state alone, which is what makes it deployable one token at a time,
and that state is built only from tokens up to $t$, so a probe never reads text
the model has not produced.

Hidden states are extracted once with the model frozen and stored, so training
reads vectors from disk. This makes it affordable to train an independent head at
\emph{every} layer in one pass and to repeat every comparison at three seeds. The
heads share no parameters and their losses are summed, which is a rearrangement
of training them separately rather than an approximation, provided the gradient
norm is clipped per head and each head keeps its own normalisation. Depth
therefore becomes an axis inside a run rather than one that multiplies the number
of runs, which matters because depth moves the results more than any other choice
studied here (Section~\ref{sec:exp-depth}).

\subsection{Which tokens the probe learns from}
\label{sec:selection}
An answer has thousands of tokens and most say nothing. Five strategies form a
ladder in which each changes exactly one decision relative to the previous one:
\texttt{all\_tokens} tiles every token; \texttt{rollout\_balanced} subsamples
negative answers and gives every answer a fixed token budget;
\texttt{random\_window} makes that budget contiguous; \texttt{frontier\_window}
anchors the window on the frontier, either centred on it or trailing it; and
\texttt{frontier\_window\_hard\_negative} additionally biases negative windows
toward spans that already look repetitive. Window widths of \TBD{64, 128, 256 and
512} are swept jointly with the rule, since an anchored window trades coverage
for earliness and its width moves the probe along that same trade-off.

\subsection{What the probe is trained toward}
\label{sec:targets}
Three label families share one contract, a per-token target and a mask marking
positions with no defined target. \textbf{Hard frontier labels} set
$y_N(t)=1 \iff f_r - t \leq N$ for a horizon $N \geq 0$, so a larger horizon asks
the probe to fire while the text still looks fine. \textbf{Soft frontier labels}
replace the step with a decay, linear or exponential over a configured length,
expressing that closer to the frontier is more degenerate without committing to a
cut. \textbf{Token-level signals} take the target from a quantity computed
independently of the frontier, such as the repetition score, which defines a
target everywhere including on healthy answers and therefore teaches a different
concept.

\subsection{An equal budget for every recipe}
The training budget is expressed in optimizer steps and target tokens per step,
and both are held identical across every recipe in a comparison. The tokens a
step sees are measured from the training stream rather than inferred from the
configured window size, because the rules contribute very different amounts per
example and deriving the budget from the configuration would quietly hand one
rule several times the budget of another. Batches are composed to a fixed
positive fraction rather than shuffled, since shuffling a population that is over
$99\%$ negative leaves most steps with no positive gradient. Class imbalance is
corrected once, in the loss weight, computed from the population the probe
actually sees after selection.

\subsection{Adapted representations}
\label{sec:adapted}
A second regime attaches low-rank adapters \citep{hu2022lora} to the layers up to
and including the probed one and trains them jointly with the head, so the
representation itself can move. An adapted run carries one probe at one depth,
since adapters below a layer change what every layer above it reads. Its control
is a frozen run at that same depth trained the same way, not one head of a
many-headed run. \TBD{the adapted runs are being repeated under the corrected
stopping rule; report only the matched comparison.}

\subsection{Which checkpoint and which depth are reported}
\label{sec:selection-rule}
A run keeps a checkpoint every \TBD{50} steps, and a rule decides which one each
depth is judged on: a depth becomes eligible once its coverage inside the loop
clears a floor, and among eligible steps the one with the best coverage of the
run-up is kept. Depth is then chosen as the eligible depth reaching the highest
such coverage. Both choices are made on validation data and reported as part of
the result.

Selecting on answer-level recall instead is a mistake we can quantify, because
that quantity saturates within the first few evaluations: it selects checkpoints
that are \TBD{6.8} times worse on coverage of the run-up, with a first alarm
\TBD{81} tokens later. This is the first of several places where the saturation
of answer-level metrics has practical consequences.

\section{Evaluation protocol}
\label{sec:protocol}

\subsection{Scores as the interface}
Evaluation never receives a model. It receives one score per token, so a probe
and a non-learned baseline go through an identical judgement by construction, and
a metric can be changed or added without recomputing a score. Scoring is
exhaustive: every token of every answer in a split, with no subsampling of
negatives and no cap, since a cap can only understate a false-alarm rate. The
protocol is checked against synthetic scorers with hand-computable answers: a
perfect oracle stepping at the frontier, oracles delayed and advanced by a fixed
number of tokens, a constant scorer, noise, and a noisy oracle that spikes on
negatives.

\subsection{Operating points}
Rather than choosing a threshold, we fix the share of \emph{healthy answers}
allowed to raise a false alarm, at $1\%$, $5\%$ and $10\%$, and solve for the
threshold that spends exactly that. Thresholds are chosen on validation and
reused unchanged everywhere else, enforced by the reporting tool rather than by
discipline. A budget framing matches the deployment question, which is how much
false-alarm cost is acceptable, and that question cannot be settled here.

A budget smaller than the share of healthy answers tied at a scorer's highest
score cannot be spent at all: the threshold is forced past the top of the range
and nothing fires anywhere, positives included. That is not a hypothetical, and
Section~\ref{sec:exp-baselines} reports which baselines it affects.

\subsection{The first alarm, and the four views}
For an answer $r$, a threshold $\tau$ and a persistence $m$, the first alarm is
\begin{equation}
  a_r(\tau,m) \;=\; \min\{\,t : p_r(t') \geq \tau \ \ \forall\, t' \in [t,\,t+m)\,\},
\end{equation}
or $\infty$ if no such $t$ exists. Four views follow.

\textbf{Detection} asks whether the alarm is finite on a degenerate answer. It
exists to confirm a scorer is not broken, not to rank scorers.

\textbf{Coverage} splits the tokens of a degenerate answer at the frontier.
\emph{In-pattern coverage} is the share of tokens at or after $f_r$ that are
flagged, the easy half and a floor to clear rather than a result. \emph{Warning
coverage} is the share flagged in a band of $128$ or $256$ tokens immediately
before $f_r$. The \emph{healthy token rate} is reported alongside both, since
flagging every token maximises them.

\textbf{Lead time} is the signed distance $a_r - f_r$, with negative meaning
early. It is reported with the share of answers that fired before the frontier
and the count that never fired, because the distribution is bimodal and a median
over it alone is misleading.

\textbf{Persistence} distinguishes a scorer that decided from one that twitched,
and is read in opposite directions on the two populations.

The reason coverage is split rather than pooled is Section~\ref{sec:onset-shape}:
the in-pattern tail dominates the positive token population and is trivially
separable, so a single blended figure reads close to perfect for almost any probe
and says nothing about the part of the problem that is hard.

\subsection{Baselines}
\label{sec:baselines}
Three structural signals computed by the labelling pipeline are scored through
the same protocol: a windowed repetition score, the longest repeated substring as
a step at the position the repeat begins, and the model's own predictive entropy,
inverted.

The windowed repetition score is \texttt{seq-rep-2} \citep{welleck2020neural}
evaluated on a moving window \citep{covington2010cutting} rather than on the whole
sequence. Its window looks \emph{ahead} of the token it labels. That is the right
choice for annotating a finished answer, where nothing is causal, and an invalid
one for a scorer a live monitor could run, because at token $t$ the tokens after
$t$ do not exist. We therefore report it in both directions and use the trailing
version in every comparison. To settle whether the shift rather than the
causality is what moves the numbers, we add \texttt{rep/l}
\citep{welleck2020neural}, the literature's per-token repetition signal, which is
backward looking by construction. \TBD{report that the two causal baselines
agree, and that the trailing windowed score is the stronger of the two and
therefore the one the probe is compared against.}

\subsection{Controls}
\label{sec:controls}
Four explanations other than an impending loop predict the same headline numbers.
Each is removed by a matched comparison rather than argued away.

\textbf{Position.} A probe's score rises with position on its own. Run-up tokens
are therefore compared against healthy tokens at the same absolute position,
grouped into bands.

\textbf{Length.} Every positive reaches the cap while healthy answers are much
shorter, so healthy answers are grouped by their own length before their
false-alarm rates are compared, with the share of each population reaching each
length reported beside the rates.

\textbf{Prompt predisposition.} Because $K$ answers share a prompt, healthy
answers whose prompt also produced a degenerate answer can be compared against
healthy answers from prompts that never looped. A probe reading imminence should
score them alike, since nothing went wrong in either.

\textbf{Domain.} Every quantity is reported per domain with its population sizes,
and any cell backed by fewer than ten positives is marked anecdotal rather than
quoted as a rate. \TBD{intervals bootstrapped over prompts rather than over
answers, since the $K$ answers to a prompt are not independent.}

\TBD{permuted-onset control: frontiers permuted among degenerate answers,
preserving the onset distribution so the control matches the real condition in
its positional statistics. This is the control task of \citet{hewitt2019designing}
for this setting, and it separates a signal about degeneration from what
supervised training extracts from any sufficiently flexible representation. The
corpus for it is built; the runs are pending.}

\section{Experiments}
\label{sec:experiments}

The programme runs in stages, each ending in a decision the next one needs. Test
splits are read only in the final stage, and no threshold and no choice of recipe
is ever made against them.

\textbf{S0, model-free signals.} What can be done with no model at all, and the
bar every later result has to clear (Section~\ref{sec:exp-baselines}).

\textbf{S1, which tokens to learn from.} The five selection strategies of
Section~\ref{sec:selection} crossed with window width, three seeds each, read on
coverage of the run-up and on depth rather than on detection
(Section~\ref{sec:exp-selection}).

\textbf{S2, what to train toward.} The horizon of the hard frontier label, the
shape and length of a soft decay, a regression target computed independently of
the frontier, and the two ways of correcting class imbalance
(Section~\ref{sec:exp-targets}).

\textbf{S3, whether adapters earn their cost.} The adapted regime of
Section~\ref{sec:adapted} against a frozen control at the same depth, trained the
same way. \TBD{in progress under the corrected stopping rule.}

\textbf{S4, the held-out test, once.} The winning configuration scored on the two
test splits with the thresholds already frozen on validation, applied unchanged.

Two further experiments sit outside the ladder. \textbf{Transfer} scores a frozen
head against other models' own corpora, both other checkpoints of the same
architecture and other model families, and trains a probe directly on those
families \TBD{in progress}. \textbf{The protocol bridge}
(Section~\ref{sec:exp-bridge}) holds one scorer fixed and turns the reporting
conventions that separate this protocol from the nearest prior work.

\section{Results}
\label{sec:results}

\subsection{Model-free signals, and which of them can be operated}
\label{sec:exp-baselines}

\begin{table}[t]
  \caption{Model-free baselines on validation at fixed per-answer false-alarm
  budgets. The forward-looking row is reported for reference only and is not a
  causal scorer. \TBD{fill; add the budget columns actually reported.}}
  \label{tab:baselines}
  \centering
  \small
  \begin{tabular}{lrrrrrrr}
    \toprule
    Scorer & AUC & Recall & Precision & In-pattern & Warning$_{256}$ & Median offset & Fired early \\
    \midrule
    windowed rep-2, forward (reference)  & \TBD{} & \TBD{} & \TBD{} & \TBD{} & \TBD{} & \TBD{} & \TBD{} \\
    windowed rep-2, trailing             & \TBD{} & \TBD{} & \TBD{} & \TBD{} & \TBD{} & \TBD{} & \TBD{} \\
    \texttt{rep/l}, $l=128$              & \TBD{} & \TBD{} & \TBD{} & \TBD{} & \TBD{} & \TBD{} & \TBD{} \\
    longest repeated substring           & \TBD{} & \TBD{} & \TBD{} & \TBD{} & \TBD{} & \TBD{} & \TBD{} \\
    inverted entropy                     & \TBD{} & \TBD{} & \TBD{} & \TBD{} & \TBD{} & \TBD{} & \TBD{} \\
    \bottomrule
  \end{tabular}
\end{table}

\TBD{state three things here: the cost of the forward window, measured; which
scorers cannot spend a $1\%$ budget and the share of healthy answers tied at
their ceiling that explains it; and the bar the probe has to clear, which is the
stronger of the two causal baselines.}

\subsection{Detection saturates; coverage of the run-up does not}
\label{sec:exp-depth}

\begin{figure}[t]
  \centering
  \figstub{4.2cm}{Depth on the horizontal axis, one panel per quantity: answer-level
  recall at a $1\%$ budget, in-pattern coverage, and warning coverage. Median and
  seed range over three seeds, each depth read at the checkpoint its own rule
  selects. The first panel is nearly flat and near one; the others move by an
  order of magnitude and peak in a broad intermediate region.}
  \caption{Answer-level detection is saturated across depth while coverage of the
  run-up is not. \TBD{caption to state the two ranges and the peak region.}}
  \label{fig:depth}
\end{figure}

\begin{table}[t]
  \caption{Selected depths for a representative frozen recipe. \TBD{fill}}
  \label{tab:depth}
  \centering
  \small
  \begin{tabular}{lrrrr}
    \toprule
    Layer & Recall & In-pattern & Warning$_{256}$ & Median offset \\
    \midrule
    \TBD{early}  & \TBD{} & \TBD{} & \TBD{} & \TBD{} \\
    \TBD{}       & \TBD{} & \TBD{} & \TBD{} & \TBD{} \\
    \TBD{middle} & \TBD{} & \TBD{} & \TBD{} & \TBD{} \\
    \TBD{}       & \TBD{} & \TBD{} & \TBD{} & \TBD{} \\
    \TBD{late}   & \TBD{} & \TBD{} & \TBD{} & \TBD{} \\
    \bottomrule
  \end{tabular}
\end{table}

\TBD{the claim is not that intermediate layers are better, which is known. It is
that the quantity this literature reports is invariant to a choice that changes
the quantity it does not report by an order of magnitude, so a paper probing the
final layer and reporting near-perfect accuracy has no way to notice.}

\subsection{Which tokens to learn from}
\label{sec:exp-selection}

\begin{table}[t]
  \caption{Token selection. Each row is a strategy at one width, averaged over
  three seeds with the spread across seeds. \TBD{fill from the paper notebook.}}
  \label{tab:selection}
  \centering
  \small
  \begin{tabular}{llrrrr}
    \toprule
    Strategy & Width & Recall & In-pattern & Warning$_{256}$ & Median offset \\
    \midrule
    \texttt{all\_tokens}                    & --      & \TBD{} & \TBD{} & \TBD{} & \TBD{} \\
    \texttt{rollout\_balanced}              & \TBD{}  & \TBD{} & \TBD{} & \TBD{} & \TBD{} \\
    \texttt{random\_window}                 & \TBD{}  & \TBD{} & \TBD{} & \TBD{} & \TBD{} \\
    \texttt{frontier\_window}, centred      & \TBD{}  & \TBD{} & \TBD{} & \TBD{} & \TBD{} \\
    \texttt{frontier\_window}, trailing     & \TBD{}  & \TBD{} & \TBD{} & \TBD{} & \TBD{} \\
    \texttt{frontier\_window\_hard\_negative} & \TBD{} & \TBD{} & \TBD{} & \TBD{} & \TBD{} \\
    \bottomrule
  \end{tabular}
\end{table}

\TBD{the ladder splits into two groups rather than five: rules that sample
broadly, and rules anchored on the frontier. The anchored rules fire earlier and
cover less, and the mechanism is that they separate degenerate from healthy
better \emph{at matched distance before the frontier}, not that they read the
loop better. State the null used: healthy answers have no frontier, so they
borrow one drawn from the frontier distribution of the degenerate answers,
without which the comparison measures score drift with position.}

\subsection{What to train toward}
\label{sec:exp-targets}

\begin{table}[t]
  \caption{Label family and horizon. \TBD{fill}}
  \label{tab:targets}
  \centering
  \small
  \begin{tabular}{llrrrr}
    \toprule
    Family & Setting & Recall & In-pattern & Warning$_{256}$ & Median offset \\
    \midrule
    hard frontier & horizon $0$        & \TBD{} & \TBD{} & \TBD{} & \TBD{} \\
    hard frontier & horizon \TBD{}     & \TBD{} & \TBD{} & \TBD{} & \TBD{} \\
    hard frontier & horizon \TBD{}     & \TBD{} & \TBD{} & \TBD{} & \TBD{} \\
    soft frontier & \TBD{decay/length} & \TBD{} & \TBD{} & \TBD{} & \TBD{} \\
    token signal  & repetition score   & \TBD{} & \TBD{} & \TBD{} & \TBD{} \\
    \bottomrule
  \end{tabular}
\end{table}

\TBD{the horizon is the one target choice that reliably moves coverage of the
run-up, and it does so monotonically in every family it was swept in, while
in-pattern coverage falls slightly and the healthy token rate does not rise. Note
explicitly that this corrects an earlier reading in which the horizon appeared to
do nothing, which came from watching the median alarm rather than the coverage.}

\subsection{What the pre-onset signal is made of}
\label{sec:exp-controls}

\begin{figure}[t]
  \centering
  \figstub{4.6cm}{One panel per control. Left: flag rate against absolute position,
  for run-up tokens and for healthy tokens, on a log scale. Middle: false-alarm
  rate by healthy-answer length band. Right: the two healthy populations of the
  prompt-sibling test as cumulative distributions.}
  \caption{Most of the apparent pre-onset signal is position, length and prompt.
  \TBD{caption to state where the residual survives and where it vanishes.}}
  \label{fig:controls}
\end{figure}

\begin{table}[t]
  \caption{Run-up tokens against healthy tokens at matched absolute position.
  \TBD{fill}}
  \label{tab:position}
  \centering
  \small
  \begin{tabular}{lrrr}
    \toprule
    Position band & Run-up flag rate & Healthy flag rate & Ratio \\
    \midrule
    \TBD{} & \TBD{} & \TBD{} & \TBD{} \\
    \TBD{} & \TBD{} & \TBD{} & \TBD{} \\
    \TBD{} & \TBD{} & \TBD{} & \TBD{} \\
    \TBD{} & \TBD{} & \TBD{} & \TBD{} \\
    \bottomrule
  \end{tabular}
\end{table}

\TBD{report the residual honestly, including the band where it disappears, and
the prompt-sibling result including the fact that adaptation does not reduce it.}

\subsection{The signal is concentrated in one kind of text}
\label{sec:exp-domain}

\begin{figure}[t]
  \centering
  \figstub{4.2cm}{Domains on the horizontal axis ordered by where their loops begin,
  one line per candidate configuration, two panels: coverage of the run-up, and
  coverage inside the loop. The two panels order the domains differently.}
  \caption{Coverage of the run-up is concentrated in one domain, which is not the
  easiest domain once inside the loop. \TBD{caption to state the share of flagged
  run-up tokens that domain supplies.}}
  \label{fig:domain}
\end{figure}

\begin{table}[t]
  \caption{Per domain, on validation and on a domain held out of training
  entirely. \TBD{fill; mark anecdotal cells.}}
  \label{tab:domain}
  \centering
  \small
  \begin{tabular}{lrrrrrr}
    \toprule
    Domain & Positives & Median onset & In-pattern & Warning$_{256}$ & Median offset & False-alarm rate \\
    \midrule
    \multicolumn{7}{l}{\emph{in domain}} \\
    \TBD{} & \TBD{} & \TBD{} & \TBD{} & \TBD{} & \TBD{} & \TBD{} \\
    \TBD{} & \TBD{} & \TBD{} & \TBD{} & \TBD{} & \TBD{} & \TBD{} \\
    \midrule
    \multicolumn{7}{l}{\emph{held out}} \\
    \TBD{} & \TBD{} & \TBD{} & \TBD{} & \TBD{} & \TBD{} & \TBD{} \\
    \bottomrule
  \end{tabular}
\end{table}

\TBD{the held-out domain is the out-of-sample test of the concentration claim,
and it was run after the claim was made on validation. State which halves of the
in-domain result replicate and which do not.}

\subsection{Converting between protocols}
\label{sec:exp-bridge}

\begin{table}[t]
  \caption{One scorer, read under successively different reporting conventions.
  No weight changes between rows. \TBD{fill}}
  \label{tab:bridge}
  \centering
  \small
  \begin{tabular}{lrrrr}
    \toprule
    & Warning$_{256}$ & Median offset & Precision & Realised FPR \\
    \midrule
    this protocol                       & \TBD{} & \TBD{} & \TBD{} & \TBD{} \\
    $+$ pooling into spans              & \TBD{} & \TBD{} & \TBD{} & \TBD{} \\
    $+$ cumulative-sum accumulation     & \TBD{} & \TBD{} & \TBD{} & \TBD{} \\
    $+$ their operating point           & \TBD{} & \TBD{} & \TBD{} & \TBD{} \\
    $+$ a balanced evaluation set       & \TBD{} & \TBD{} & \TBD{} & \TBD{} \\
    \bottomrule
  \end{tabular}
\end{table}

The accumulation rung uses the one-sided cumulative-sum recursion
\citep{page1954continuous},
\begin{equation}
  S_i \;=\; \max\!\left(0,\; S_{i-1} + (x_i - k)\right),\qquad S_0 = 0,
\end{equation}
where $x_i$ is the mean probe score over span $i$ and $k$ is a reference level
measured as the median pooled score over healthy answers. The firing threshold is
solved for the false-alarm budget, as for every other row, which is what keeps
the rows comparable.

\TBD{state the ordering of the terms and, most importantly, that the precision
column is where the cost of a loose operating point appears and that a balanced
evaluation set removes it from view. Add the observation that a balanced set
thinned to the number of positives cannot resolve a strict budget at all.}

This measures one scorer under other reporting choices. It is not an evaluation
of another detector and says nothing about what another model would do on this
data.

\subsection{Transfer}
\label{sec:exp-transfer}

\begin{table}[t]
  \caption{A frozen head scored against other models' own corpora, and probes
  trained directly on those corpora. \TBD{fill; the cross-family rows are
  pending a rebuild.}}
  \label{tab:transfer}
  \centering
  \small
  \begin{tabular}{llrrr}
    \toprule
    Target & Relation & Recall & In-pattern & Warning$_{256}$ \\
    \midrule
    \TBD{} & same architecture, other weights & \TBD{} & \TBD{} & \TBD{} \\
    \TBD{} & other family, transferred head   & \TBD{} & \TBD{} & \TBD{} \\
    \TBD{} & other family, probe trained on it & \TBD{} & \TBD{} & \TBD{} \\
    \bottomrule
  \end{tabular}
\end{table}

\TBD{keep the two questions apart: whether a learned direction survives a change
of weights, and whether the phenomenon is readable at all in another family. Only
the second is answered by training on that family.}

\section{Limitations}
\label{sec:limitations}

\TBD{to be written, covering at least: one model family for the main results; the
number of positives behind the domain claim; the frontier as a judge-produced
reference point without human validation; the scope of the task as persistent
degeneration that reaches the cap; and the fact that the selection of depth and
checkpoint is made on the same validation split as everything else.}

\section{Conclusion}
\TBD{written last.}

\begin{ack}
\TBD{withheld for submission.}
\end{ack}

\bibliographystyle{plainnat}
\bibliography{references}

\appendix

\section{Structural repetition rules, measured}
\TBD{the deployed-rules result: detection against the judge, firing rate on
naturally terminated text, and firing rate on capped answers the judge confirms
are not degenerate. State that a disagreement bounds the joint error of rule and
judge, and that these numbers come from a different protocol than the probe
results and must not be placed in one table with them.}

\section{Corpus documentation}
\TBD{why an end-of-sequence answer is treated as a negative, why only capped
answers are judged, judge reliability and refusals, and the three-build
comparison.}

\section{Additional results}
\TBD{persistence, all window widths, soft decay shapes, every depth and
checkpoint, per-domain intervals, and the adapted regime in full.}

\end{document}
